% ============================================================
%  CAPITOLO 4 --- PROGETTAZIONE LOGICA
% ============================================================
\chapter{Progettazione logica}
\setlength{\headheight}{25.5pt}
\raggedbottom

% -- 4.1 Stima del volume dei dati ---------------------------
\section{Stima del volume dei dati}

La tabella seguente riporta la stima del numero di istanze per ogni entit\`{a} e
associazione, basata su un'ipotetica stagione balneare di medio-grandi dimensioni
(~200 ombrelloni, ~1000 clienti registrati).

\begin{table}[H]
  \centering
  \caption{Stima del volume dei dati}
  \label{tab:volumi}
  \rowcolors{2}{rowalt}{white}
  \begin{tabularx}{\textwidth}{L{4cm} C{2cm} R{3cm} L{5cm}}
    \toprule
    \rowcolor{primary!15}
    \textbf{Concetto} & \textbf{Tipo} & \textbf{Volume} & \textbf{Note} \\
    \midrule
    \entita{Proprietario}    & E & 2--5      & Soci dello stabilimento \\
    \entita{Cliente}         & E & 1.000     & Clienti registrati a stagione \\
    \entita{Ombrellone}      & E & 200       & Postazioni totali \\
    \entita{Prenotazione}    & E & 3.000     & ~15 prenotazioni/ombrellone \\
    \entita{Abbonamento}     & E & 150       & ~15\% dei clienti \\
    \entita{Dipendente}      & E & 20        & Personale stagionale \\
    \entita{Turno}           & E & 1.800     & ~90 turni/dipendente/stagione \\
    \entita{Prodotto}        & E & 80        & Articoli a catalogo \\
    \entita{Fornitore}       & E & 10        & Fornitori abituali \\
    \entita{Ordine}          & E & 100       & ~2 ordini/settimana \\
    \relazione{Contiene}     & A & 300       & ~3 prodotti/ordine \\
    \bottomrule
  \end{tabularx}
\end{table}

% -- 4.2 Operazioni principali -------------------------------
\section{Descrizione delle operazioni principali e stima della loro frequenza}

\begin{table}[H]
  \centering
  \caption{Operazioni principali e frequenza}
  \label{tab:operazioni}
  \rowcolors{2}{rowalt}{white}
  \begin{tabularx}{\textwidth}{C{0.6cm} L{6cm} C{2.5cm} C{2.5cm}}
    \toprule
    \rowcolor{primary!15}
    \textbf{N.} & \textbf{Operazione} & \textbf{Frequenza} & \textbf{Tipo} \\
    \midrule
    O1 & Inserimento nuova prenotazione            & 30/giorno   & Scrittura \\
    O2 & Verifica disponibilit\`{a} ombrelloni per data& 50/giorno   & Lettura \\
    O3 & Visualizzazione dipendenti in turno per data & 5/giorno & Lettura \\
    O4 & Inserimento/modifica turno dipendente     & 10/giorno   & Scrittura \\
    O5 & Controllo giacenze magazzino              & 3/giorno    & Lettura \\
    O6 & Inserimento nuovo ordine                  & 2/settimana & Scrittura \\
    O7 & Inserimento nuovo cliente                 & 10/giorno   & Scrittura \\
    O8 & Visualizzazione storico prenotazioni cliente & 5/giorno & Lettura \\
    \bottomrule
  \end{tabularx}
\end{table}

% -- 4.3 Schemi di navigazione -------------------------------
\section{Schemi di navigazione e tabelle degli accessi}

\subsection{O2 --- Verifica disponibilit\`{a} ombrelloni per una data}

Questa operazione verifica quali ombrelloni non hanno prenotazioni attive
in una determinata data.

\textbf{Schema di navigazione:} si accede all'entit\`{a} \entita{Ombrellone},
poi all'associazione \relazione{PrenotaOmb} e infine all'entit\`{a}
\entita{Prenotazione} per verificare le date.

\begin{table}[H]
  \centering
  \caption{Tabella degli accessi --- O2: Disponibilit\`{a} ombrelloni}
  \label{tab:accessi-o2}
  \rowcolors{2}{rowalt}{white}
  \begin{tabularx}{\textwidth}{L{4cm} C{2cm} C{2cm} C{2cm} C{2cm}}
    \toprule
    \rowcolor{primary!15}
    \textbf{Concetto} & \textbf{Tipo} & \textbf{L/S} & \textbf{Accessi} & \textbf{Tot/mese} \\
    \midrule
    \entita{Ombrellone}   & E & L & 200    & 300.000 \\
    \entita{Prenotazione} & E & L & 3.000  & 4.500.000 \\
    \relazione{PrenotaOmb}& A & L & 3.000  & 4.500.000 \\
    \bottomrule
  \end{tabularx}
\end{table}

\subsection{O3 --- Dipendenti in turno per una data specifica}

\begin{table}[H]
  \centering
  \caption{Tabella degli accessi --- O3: Dipendenti in turno}
  \label{tab:accessi-o3}
  \rowcolors{2}{rowalt}{white}
  \begin{tabularx}{\textwidth}{L{4cm} C{2cm} C{2cm} C{2cm} C{2cm}}
    \toprule
    \rowcolor{primary!15}
    \textbf{Concetto} & \textbf{Tipo} & \textbf{L/S} & \textbf{Accessi} & \textbf{Tot/mese} \\
    \midrule
    \entita{Turno}      & E & L & 1.800  & 27.000 \\
    \entita{Dipendente} & E & L & 20     & 300 \\
    \relazione{LavoraIn}& A & L & 1.800  & 27.000 \\
    \bottomrule
  \end{tabularx}
\end{table}

% -- 4.4 Analisi delle ridondanze ----------------------------
\section{Analisi delle ridondanze}

\textbf{Ridondanza candidata:} l'attributo \attributo{prezzoTotale} in
\entita{Prenotazione} \`{e} calcolabile come:
\[
  \text{prezzoTotale} = \text{tariffaGiornaliera}(\text{posizione}) \times
  (\text{dataFine} - \text{dataInizio})
\]

\textbf{Decisione:} si mantiene l'attributo ridondante \attributo{prezzoTotale}
perch\'e la tariffa potrebbe variare nel tempo e si vuole registrare il prezzo
concordato al momento della prenotazione. Questo evita problemi di consistenza
se le tariffe cambiano.

% -- 4.5 Raffinamento schema ---------------------------------
\section{Raffinamento dello schema}

\subsection{Eliminazione degli identificatori esterni}
Nessuna entit\`{a} nello schema utilizza identificatori esterni: tutte le entit\`{a}
sono identificate da una chiave surrogata (\texttt{id} intero auto-incrementale).

\subsection{Attributi composti}
L'attributo composto \textit{posizione} di \entita{Ombrellone} viene decomposto
in \attributo{fila} (CHAR) e \attributo{numero} (INTEGER).

\subsection{Gerarchie}
Non sono presenti gerarchie di generalizzazione nello schema.

\subsection{Scelta delle chiavi}
Ogni entit\`{a} viene identificata da una chiave surrogata \texttt{id SERIAL PRIMARY KEY}.
Le chiavi naturali vengono mantenute come vincoli UNIQUE dove applicabile
(es. coppia \texttt{fila + numero} per \entita{Ombrellone}, \texttt{email}
per \entita{Cliente} e \entita{Proprietario}).

% -- 4.6 Traduzione in relazioni -----------------------------
\section{Traduzione di entit\`{a} e associazioni in relazioni}

\begin{longtable}{L{3.5cm} L{10cm}}
  \caption{Traduzione dello schema E/R in relazioni}
  \label{tab:traduzione} \\
  \toprule
  \rowcolor{primary!15}
  \textbf{Relazione} & \textbf{Schema (PK sottolineata, FK in corsivo)} \\
  \midrule
  \endfirsthead
  \toprule
  \rowcolor{primary!15}
  \textbf{Relazione} & \textbf{Schema} \\
  \midrule
  \endhead
  \bottomrule
  \endfoot

  PROPRIETARIO   & (\chiave{id}, nome, cognome, email, passwordHash) \\[4pt]
  CLIENTE        & (\chiave{id}, nome, cognome, telefono, email) \\[4pt]
  OMBRELLONE     & (\chiave{id}, fila, numero, numeroLettini, numeroSedie) \\[4pt]
  PRENOTAZIONE   & (\chiave{id}, \chiaveest{idCliente}, \chiaveest{idOmbrellone},
                   dataInizio, dataFine, prezzoTotale) \\[4pt]
  ABBONAMENTO    & (\chiave{id}, \chiaveest{idCliente}, \chiaveest{idOmbrellone},
                   dataInizio, dataFine, prezzo, tipo) \\[4pt]
  DIPENDENTE     & (\chiave{id}, nome, cognome, codiceFiscale, telefono, email) \\[4pt]
  TURNO          & (\chiave{id}, \chiaveest{idDipendente}, data,
                   orarioInizio, orarioFine) \\[4pt]
  PRODOTTO       & (\chiave{id}, nome, descrizione, quantita, prezzoUnitario,
                   sogliaMinimaRiordino) \\[4pt]
  FORNITORE      & (\chiave{id}, ragioneSociale, telefono, email, indirizzo) \\[4pt]
  ORDINE         & (\chiave{id}, \chiaveest{idFornitore}, \chiaveest{idProprietario},
                   dataOrdine, statoConsegna) \\[4pt]
  CONTIENE       & (\chiave{\chiaveest{idOrdine}}, \chiave{\chiaveest{idProdotto}},
                   quantitaOrdinata, prezzoUnitarioAlMomento) \\[4pt]
\end{longtable}

% -- 4.7 Schema relazionale finale ---------------------------
\section{Schema relazionale finale}

\begin{figure}[H]
  \centering
  % \includegraphics[width=\textwidth]{immagini/schema_relazionale.png}
  \fbox{\parbox{0.95\textwidth}{\centering\vspace{3cm}
    \textit{[Inserire immagine: schema\_relazionale.png]}\\
    Diagramma dello schema relazionale con chiavi e foreign key
  \vspace{3cm}}}
  \caption{Schema relazionale finale}
  \label{fig:schema-relazionale}
\end{figure}

% -- 4.8 SQL: CREATE TABLE ------------------------------------
\section{Costruzione delle tabelle in SQL}

\begin{lstlisting}[caption={Creazione delle tabelle principali},
                   label={lst:create-tables}]
-- Proprietario
CREATE TABLE Proprietario (
    id            SERIAL PRIMARY KEY,
    nome          VARCHAR(50)  NOT NULL,
    cognome       VARCHAR(50)  NOT NULL,
    email         VARCHAR(100) NOT NULL UNIQUE,
    passwordHash  VARCHAR(255) NOT NULL
);

-- Cliente
CREATE TABLE Cliente (
    id        SERIAL PRIMARY KEY,
    nome      VARCHAR(50)  NOT NULL,
    cognome   VARCHAR(50)  NOT NULL,
    telefono  VARCHAR(20),
    email     VARCHAR(100) UNIQUE
);

-- Ombrellone
CREATE TABLE Ombrellone (
    id             SERIAL PRIMARY KEY,
    fila           CHAR(1)  NOT NULL,
    numero         INTEGER  NOT NULL,
    numeroLettini  INTEGER  NOT NULL DEFAULT 2,
    numeroSedie    INTEGER  NOT NULL DEFAULT 2,
    UNIQUE (fila, numero)
);

-- Prenotazione
CREATE TABLE Prenotazione (
    id              SERIAL PRIMARY KEY,
    idCliente       INTEGER NOT NULL
                    REFERENCES Cliente(id) ON DELETE RESTRICT,
    idOmbrellone    INTEGER NOT NULL
                    REFERENCES Ombrellone(id) ON DELETE RESTRICT,
    dataInizio      DATE    NOT NULL,
    dataFine        DATE    NOT NULL,
    prezzoTotale    DECIMAL(8,2),
    CONSTRAINT chk_date CHECK (dataFine >= dataInizio)
);

-- Dipendente
CREATE TABLE Dipendente (
    id            SERIAL PRIMARY KEY,
    nome          VARCHAR(50)  NOT NULL,
    cognome       VARCHAR(50)  NOT NULL,
    codiceFiscale CHAR(16)     UNIQUE,
    telefono      VARCHAR(20),
    email         VARCHAR(100)
);

-- Turno
CREATE TABLE Turno (
    id            SERIAL PRIMARY KEY,
    idDipendente  INTEGER NOT NULL
                  REFERENCES Dipendente(id) ON DELETE CASCADE,
    data          DATE    NOT NULL,
    orarioInizio  TIME    NOT NULL,
    orarioFine    TIME    NOT NULL,
    CONSTRAINT chk_orari CHECK (orarioFine > orarioInizio)
);

-- Prodotto
CREATE TABLE Prodotto (
    id                    SERIAL PRIMARY KEY,
    nome                  VARCHAR(100) NOT NULL,
    descrizione           TEXT,
    quantita              INTEGER NOT NULL DEFAULT 0,
    prezzoUnitario        DECIMAL(8,2),
    sogliaMinimaRiordino  INTEGER DEFAULT 5
);

-- Fornitore
CREATE TABLE Fornitore (
    id             SERIAL PRIMARY KEY,
    ragioneSociale VARCHAR(150) NOT NULL,
    telefono       VARCHAR(20),
    email          VARCHAR(100),
    indirizzo      TEXT
);

-- Ordine
CREATE TABLE Ordine (
    id               SERIAL PRIMARY KEY,
    idFornitore      INTEGER NOT NULL
                     REFERENCES Fornitore(id) ON DELETE RESTRICT,
    idProprietario   INTEGER NOT NULL
                     REFERENCES Proprietario(id) ON DELETE RESTRICT,
    dataOrdine       DATE    NOT NULL DEFAULT CURRENT_DATE,
    statoConsegna    VARCHAR(20) NOT NULL DEFAULT 'in_attesa'
                     CHECK (statoConsegna IN ('in_attesa','consegnato','annullato'))
);

-- Contiene (associazione Ordine--Prodotto)
CREATE TABLE Contiene (
    idOrdine               INTEGER NOT NULL
                           REFERENCES Ordine(id) ON DELETE CASCADE,
    idProdotto             INTEGER NOT NULL
                           REFERENCES Prodotto(id) ON DELETE RESTRICT,
    quantitaOrdinata       INTEGER NOT NULL CHECK (quantitaOrdinata > 0),
    prezzoUnitarioAlMomento DECIMAL(8,2),
    PRIMARY KEY (idOrdine, idProdotto)
);
\end{lstlisting}

% -- 4.9 Query SQL -------------------------------------------
\section{Traduzione delle operazioni in query SQL}

\subsection{O2 --- Ombrelloni disponibili in un intervallo di date}

\begin{lstlisting}[caption={O2: Ombrelloni disponibili per una data}]
-- Parametri: :data_inizio, :data_fine
SELECT o.id, o.fila, o.numero, o.numeroLettini, o.numeroSedie
FROM   Ombrellone o
WHERE  o.id NOT IN (
    SELECT p.idOmbrellone
    FROM   Prenotazione p
    WHERE  p.dataInizio <= :data_fine
    AND    p.dataFine   >= :data_inizio
)
ORDER BY o.fila, o.numero;
\end{lstlisting}

\subsection{O3 --- Dipendenti in turno per una data specifica}

\begin{lstlisting}[caption={O3: Dipendenti in turno per una data}]
-- Parametro: :data
SELECT d.nome, d.cognome, t.orarioInizio, t.orarioFine
FROM   Dipendente d
JOIN   Turno t ON t.idDipendente = d.id
WHERE  t.data = :data
ORDER BY t.orarioInizio;
\end{lstlisting}

\subsection{O5 --- Prodotti sotto la soglia minima in magazzino}

\begin{lstlisting}[caption={O5: Mancanze in magazzino}]
SELECT p.nome, p.quantita, p.sogliaMinimaRiordino,
       (p.sogliaMinimaRiordino - p.quantita) AS quantita_mancante
FROM   Prodotto p
WHERE  p.quantita < p.sogliaMinimaRiordino
ORDER BY quantita_mancante DESC;
\end{lstlisting}

\subsection{O1 --- Inserimento nuova prenotazione}

\begin{lstlisting}[caption={O1: Inserimento prenotazione con calcolo prezzo}]
-- Inserisce una nuova prenotazione
-- Il prezzo totale viene passato dall'applicazione dopo il calcolo
INSERT INTO Prenotazione (idCliente, idOmbrellone,
                          dataInizio, dataFine, prezzoTotale)
VALUES (:id_cliente, :id_ombrellone,
        :data_inizio, :data_fine, :prezzo_totale);
\end{lstlisting}

\subsection{O8 --- Storico prenotazioni di un cliente}

\begin{lstlisting}[caption={O8: Storico prenotazioni di un cliente}]
-- Parametro: :id_cliente
SELECT o.fila, o.numero, p.dataInizio, p.dataFine, p.prezzoTotale
FROM   Prenotazione p
JOIN   Ombrellone o ON o.id = p.idOmbrellone
WHERE  p.idCliente = :id_cliente
ORDER BY p.dataInizio DESC;
\end{lstlisting}
